% easychair.tex,v 3.5 2017/03/15

\documentclass{easychair}

\usepackage{hyperref}
\usepackage{booktabs}
\usepackage{footnoterange}
\usepackage{todonotes}
\usepackage{cite}
\usepackage{siunitx}
\usepackage{amsmath,amssymb}
\usepackage[T1]{fontenc}
\usepackage[utf8]{inputenc}

\title{The World Model Calculus: \\ Autoformalizing and Extending Probabilistic Logic Networks
\thanks{This document has been largely generated by AI under Zar's direction and will need to be checked carefully before camera-ready versions.  I believe the core work in progress is worth sharing.  There is a lengthy exposition in the form of a \href{https://github.com/zariuq/ai-agents/blob/main/lean-projects/mettapedia/papers/wm-pln-book_v3.pdf}{book draft} online..}
}

\author{
Zarathustra Amadeus Goertzel
\and Oruži\thanks{A collaboration of AI coworkers consisting of at least Claude 4.x, Codex 5.2+, and ChatGPT Pro.}
}

\institute{
  Czech Technical University in Prague, Czech Republic
 }

\authorrunning{Goertzel \& Oruži}
\titlerunning{World Model Calculus}

\newcommand{\WM}{\mathsf{WM}}
\newcommand{\PLN}{\mathsf{PLN}}
\newcommand{\Ev}{\mathcal{E}}
\newcommand{\State}{\mathcal{S}}
\newcommand{\Query}{\mathcal{Q}}
\newcommand{\code}[1]{\texttt{#1}}
\sloppy

\begin{document}
\maketitle

\vspace{-3mm}
\paragraph*{Abstract:}
We propose a World Model (WM) calculus that is able to cover other approaches to reasoning over uncertainty, given the correct world model.  The WM calculus originated from an attempt to autoformalize Probabilistic Logic Networks (PLN) in Lean~4, yet also seems capable of representing Markov Logic Networks and ProbLog.  PLN inference rules are viewed as conditioned queries to world models: the core logic is one of evidence revision and querying.  We hope that the WM calculus helps to advance the state of reasoning over uncertainty, as well as providing clear no-go theorems as to when local chaining is insufficient and exploring how higher-order probabilities may guide the integration of machine learning predictors into reasoning systems.

\vspace{-3mm}
\paragraph*{Introduction:}
This abstract presents what I think is the right core semantics for WM--PLN if we want to use PLN seriously inside theorem proving, scientific reasoning, and AGI-style inference control.  The old PLN strength/confidence truth value was a good engineering instinct: a claim supported by $4$ positive examples out of $5$ should not behave the same as one supported by $400$ out of $500$.  The trouble is that it is too easy to treat a local truth-value pair as if it were the semantics itself.  Here the primitive object is instead a revisable world-model state $W$, a query $q$, and an extracted evidence object $\mathrm{evidence}(W,q) \in \Ev$.  Strength, confidence, intervals, and probabilities are then views of evidence rather than the semantic object itself.  The current Lean development centers on modules such as \code{PLNWorldModelGeneric.lean}, \code{PLNBayesNetFastRules.lean}, \code{ProbLogDistributionSemantics.lean}, \code{MarkovLogicClauseWorldModel.lean}, and \code{MarkovLogicInfiniteWorldModel.lean}.

\vspace{-3mm}
\paragraph*{Core WM calculus:}
In the main additive regime, revision combines states and evidence extraction commutes with revision:
\[
  \mathrm{evidence}(W_1 + W_2,q)
  = \mathrm{evidence}(W_1,q) + \mathrm{evidence}(W_2,q).
\]
This is the semantic contract behind the \code{BinaryWorldModel}/\code{PLNWorldModelGeneric} interface.  It makes explicit something that was only implicit in traditional PLN: the truth value should come \emph{after} query answering, not before it.  In particular, binary Beta-style evidence, Dirichlet-style evidence, Normal--Gamma style evidence, provenance annotations, and source-trust layers can all live in the evidence carrier, while the familiar PLN strength and confidence are recovered as views.  This gives a cleaner semantic home to the original PLN intuition that probability and amount-of-support should not be conflated~\cite{plnbook,knuthskilling}.

\vspace{-3mm}
\paragraph*{Certified chaining:}
A key theorem is that there is no local complete link-deduction rule: there is no function of only the local evidence for $A$, $B$, $C$, $A \rightarrow B$, and $B \rightarrow C$ that returns the exact evidence for $A \rightarrow C$ for all world models.  Correlations live in the joint state, not in five local summaries.  This gives what I think is the clean answer to a longstanding PLN question: chaining works when it is a theorem about a model class.  So the right status of a PLN rule is not ``universal truth-value algebra'' but a $\Sigma$-guarded rewrite.  For Bayesian-network chains and forks, the classical PLN deduction formula becomes exact when screening-off and positivity conditions hold, as in \code{chainBN\_plnDeductionStrength\_exact} and \code{forkBN\_plnDeductionStrength\_exact}.  For collider patterns, hidden confounders, or untracked overlap, the system should not pretend exactness: it should either widen to an interval, keep an error ledger, or abstain~\cite{pearl,halpern}.

\vspace{-3mm}
\paragraph*{Bridges to ProbLog and Markov logic:}
This is also how I think ProbLog and MLNs fit cleanly into the picture.  ProbLog's distribution semantics gives a world-model class where WM query strength agrees with ProbLog query probability, with the BDD / weighted-model-counting layer giving the compiled execution path~\cite{problog}.  Markov logic becomes another world-model class: finite clause-mass semantics on one side and infinite Gibbs/DLR semantics on the other, both packaged into the same evidence/query interface~\cite{mln}.  So WM--PLN is not trying to compete with ProbLog or MLN as yet another probabilistic logic.  The point is rather to give them a common semantic address space in which PLN-style fast rules can be proved, constrained, or refused.

\vspace{-3mm}
\paragraph*{Higher-order chaining and AITP relevance:}
The next problem, and the one I especially want to discuss at AITP, is how to most correctly do higher-order chaining of probabilistic inference.  In practice we want uncertainties not only about atoms, but also about rule applicability, conjecture plausibility, premise usefulness, and which proof-search policy to fire next.  My current view is that this should not be collapsed into one mysterious confidence number.  Higher-order uncertainty should live over admissibility regimes, hidden context, and trust in the guide; every derived edge should carry its side conditions, provenance, and mode of use.  Then large uncertain steps can guide search, while the small kernel still checks what it can actually justify.  This also seems to me the right meeting point between the old PLN agenda and the AITP agenda of making learned or heuristic guidance answerable to formal structure rather than merely helpful by accident~\cite{malarea,enigma}.

\vspace{-3mm}
\paragraph*{Futureward:}
This is deliberately a work-in-progress / project proposal.  The main technical question is which probabilistic shortcuts a practical ATP should try to certify, and how much structure a premise-selection or chaining system must export so that its uncertainty estimates can become proof-carrying rather than merely heuristic.  My hope is that WM--PLN can give a mathematically cleaner home for PLN while still preserving the engineering spirit that made PLN useful in the first place.

\vspace{-3mm}
\paragraph*{Acknowledgements}
Many thanks to Ben Goertzel for the original PLN vision, to Nil Geisweiller for ongoing PLN inference-control work, and to Josef Urban for many related conversations around ATP and learned guidance.
Thanks to OpenAI for providing amazing AI services at reasonable costs.
This work is supported by Dar Novamente --- Cognitive Inference Control.

\newpage
\bibliographystyle{plain}
\begin{thebibliography}{9}
\bibitem{plnbook}
B.~Goertzel, M.~Ikl\'{e}, I.~F. Goertzel, and A.~Heljakka.
\newblock \emph{Probabilistic Logic Networks: A Comprehensive Framework for Uncertain Inference}.
\newblock Springer, 2008.

\bibitem{knuthskilling}
K.~H. Knuth and J.~Skilling.
\newblock Foundations of inference.
\newblock \emph{Axioms}, 1(1):38--73, 2012.

\bibitem{pearl}
J.~Pearl.
\newblock \emph{Probabilistic Reasoning in Intelligent Systems}.
\newblock Morgan Kaufmann, 1988.

\bibitem{halpern}
J.~Y. Halpern.
\newblock \emph{Reasoning about Uncertainty}.
\newblock MIT Press, 2003.

\bibitem{problog}
L.~De Raedt, A.~Kimmig, and H.~Toivonen.
\newblock ProbLog: A probabilistic Prolog and its application in link discovery.
\newblock In \emph{IJCAI}, pages 2468--2473, 2007.

\bibitem{mln}
M.~Richardson and P.~Domingos.
\newblock Markov logic networks.
\newblock \emph{Machine Learning}, 62(1--2):107--136, 2006.

\bibitem{malarea}
J.~Urban, G.~Sutcliffe, P.~Pudl\'{a}k, and J.~Vysko\v{c}il.
\newblock MaLARea SG1: Machine learner for automated reasoning with semantic guidance.
\newblock In \emph{IJCAR}, pages 441--456, 2008.

\bibitem{enigma}
J.~Jakub\r{u}v and J.~Urban.
\newblock ENIGMA: Efficient learning-based inference guiding machine.
\newblock In \emph{CICM}, pages 292--302, 2017.
\end{thebibliography}

\end{document}
